% Independent proof fragment. 2026-09-19.
% Input: 04_FULL_METRIC_POSITIVE_TRACE_MINIMUM.md, SHA256
% 7d26e6933a65c7e27a739966560bdb07e2b83b344f926b743137621534203d13.
% This fragment uses amsmath, amssymb and no nonstandard macros.

\section{Every minimizing full-packet metric, exact excess, and a rank-two minimizer}
\label{sec:pte}

All matrices and maps below act on the complete original packet.  Fix
\(k=4\ell+1\geq1\), \(m_0\geq1\), \(\delta>0\), and \(\gamma>0\), and put
\begin{equation}\tag{PTE1}\label{eq:pte1}
 e=1+k(m_0-1),\qquad q=e(k+1)^2,\qquad d=q/2,
 \quad \lambda_{a,b}=\frac{k}{2}+(2a-k)\delta+i(2b-k)\gamma,
 \quad 0\leq a,b\leq k.
\end{equation}
Let \(\chi(S)=\prod_{a,b}(S-\lambda_{a,b})^e\),
\(E=\mathbb C[S]/(\chi)\), and let \(A:E\longrightarrow E\) be multiplication
by \(S\).  The fixed original polynomial basis is used whenever a metric
matrix is written without an explicit change of frame.  Set
\begin{equation}\tag{PTE2}\label{eq:pte2}
 B=A-\frac{k}{2}I,
 \quad H_G=G^{-1/2}(B^*G+GB)G^{-1/2},
 \quad L=\frac{\delta q(k+1)}2,
 \quad \mathcal E(G)=\operatorname{Tr}(H_G)_+-L
 \qquad(G>0).
\end{equation}
Here \(T_+\) and \(T_-\), for a Hermitian matrix \(T\), denote the positive
semidefinite matrices obtained by replacing its eigenvalues \(t\) by
\(\max(t,0)\) and \(\max(-t,0)\), respectively.  Thus \(T=T_+-T_-\).
The symbol \(\mathcal E(G)\) is an excess in the positive trace, and is
unrelated to the packet space \(E\).

Let \(E_+\) and \(E_-\) be the sums of all complete primary spaces with
\(a>k/2\) and \(a<k/2\), respectively.  Then
\(E=E_+\oplus E_-\), with both dimensions equal to \(d\).  Write
\(B_+=B|_{E_+}\), \(B_-=B|_{E_-}\), and \(C_-=-B_-\).
Each eigenvalue of \(B_+\) and \(C_-\) has positive real part.  In
particular their smallest real parts are \(\delta\).  A sum without
truncating a primary block gives
\begin{equation}\tag{PTE3}\label{eq:pte3}
 2\operatorname{Re}\operatorname{Tr}B_+
 =2e(k+1)\delta\sum_{a=(k+1)/2}^{k}(2a-k)
 =\frac{\delta e(k+1)^3}{2}=L,
 \quad 2\operatorname{Re}\operatorname{Tr}B_-=-L,
 \quad \operatorname{Tr}H_G=0.
\end{equation}
Indeed, the positive odd integers in the sum are \(1,3,\ldots,k\), whose
sum is \(((k+1)/2)^2\).  The negative sum is its negative.  The nilpotent
parts have zero trace; this assertion keeps their full multiplicities.

\subsection{An exact metric frame and the excess identity}

Choose any fixed bases of the original complete spaces \(E_+\) and
\(E_-\).  In their direct sum frame the metric and the operator are
\begin{equation}\tag{PTE4}\label{eq:pte4}
 G=\begin{pmatrix}P&X\\X^*&Y\end{pmatrix}>0,
 \qquad B=\begin{pmatrix}B_+&0\\0&B_-\end{pmatrix},
 \qquad K=P^{-1}X,
 \quad S=Y-X^*P^{-1}X>0.
\end{equation}
Positivity of \(P\) follows by restriction.  For each \(v\ne0\), the
vector \((-Kv,v)\) has squared metric norm \(v^*Sv>0\), proving
positivity of \(S\).  Define the exact map from an ordinary coordinate
sum to this original frame by
\begin{equation}\tag{PTE5}\label{eq:pte5}
 F=\begin{pmatrix}P^{-1/2}&-KS^{-1/2}\\0&S^{-1/2}\end{pmatrix}.
 \qquad F^*GF=I.
\end{equation}
The first summand is sent onto \(E_+\), and the second onto its actual
\(G\)-orthogonal complement.  Direct multiplication gives
\begin{equation}\tag{PTE6}\label{eq:pte6}
 F^{-1}BF=
 \begin{pmatrix}\widehat B_+&C\\0&\widehat B_-\end{pmatrix},
 \quad \widehat B_+=P^{1/2}B_+P^{-1/2},
 \quad \widehat B_-=S^{1/2}B_-S^{-1/2},
 \quad C=P^{1/2}(KB_--B_+K)S^{-1/2}.
\end{equation}
Consequently
\begin{equation}\tag{PTE7}\label{eq:pte7}
 H=F^*(B^*G+GB)F
 =\begin{pmatrix}D_+&C\\C^*&D_-\end{pmatrix},
 \quad D_\pm=\widehat B_\pm+\widehat B_\pm^*,
 \quad \operatorname{Tr}D_+=L,\quad \operatorname{Tr}D_-=-L.
\end{equation}
The matrix \(U=G^{1/2}F\) is unitary and \(H=U^*H_GU\), so no spectral
quantity has changed.  If the direct sum frame in (PTE4) is not the
original polynomial frame, its fixed basis-change matrix is inserted
on both sides of \(F\); the same identity follows.

Let \(P_0=\operatorname{diag}(I_d,0)\), \(Q_0=I-P_0\), and
\(J_0=P_0-Q_0\).  The following identities and inequalities are exact:
\begin{align}
 \mathcal E(G)
 &=\operatorname{Tr}(Q_0H_+)+\operatorname{Tr}(P_0H_-)
 =\|H_+^{1/2}Q_0\|_{\mathrm F}^2
   +\|H_-^{1/2}P_0\|_{\mathrm F}^2\geq0,
 \tag{PTE8}\label{eq:pte8}\\
 \mathcal E(G)
 &=\frac14\| (\operatorname{sgn}H-J_0)|H|^{1/2}\|_{\mathrm F}^2.
 \tag{PTE9}\label{eq:pte9}
\end{align}
In (PTE9), \(\operatorname{sgn}(0)=0\), and \(\|\cdot\|_{\mathrm F}\)
is the Frobenius norm.  To prove (PTE8), subtract
\(L=\operatorname{Tr}(P_0H_+)-\operatorname{Tr}(P_0H_-)\) from
\(\operatorname{Tr}H_+\).  Each trace on its right is nonnegative, because
it is the displayed squared norm.  For (PTE9), expansion of the squared
norm yields
\[
 2\operatorname{Tr}|H|-2\operatorname{Tr}(J_0H).
\]
The zero spectral subspace contributes zero even though
\((\operatorname{sgn}H)^2\) is not the identity on that subspace.
Now \(\operatorname{Tr}|H|=2\operatorname{Tr}H_+\), because
\(\operatorname{Tr}H=0\), and \(\operatorname{Tr}(J_0H)=2L\).
The expansion is therefore \(4\mathcal E(G)\), as asserted.

\subsection{Equality, including zero eigenvalues}

The exact set of minimizing metrics is
\begin{equation}\tag{PTE10}\label{eq:pte10}
 G(E_+,E_-)=0,
 \qquad B_+^*G_++G_+B_+\succeq0,
 \qquad C_-^*G_-+G_-C_-\succeq0,
 \qquad G_+>0, G_->0.
\end{equation}
Here \(G_\pm\) are the restrictions in the same fixed primary-space
bases; the semidefinite signs are deliberately non-strict.

First suppose \(\mathcal E(G)=0\).  By (PTE8),
\(H_+^{1/2}Q_0=0\) and \(H_-^{1/2}P_0=0\).  Taking adjoints also
annihilates the opposite sides.  Thus \(H_+\) is supported in \(P_0\),
\(H_-\) is supported in \(Q_0\), and
\(C=P_0HQ_0=0\).  The Sylvester map
\(K\mapsto B_+K-KB_-\) is invertible: a fully evaluated inverse is
proved in (PTE23) below.  Equation (PTE6) therefore gives \(K=0\), hence
\(X=0\).  Equations (PTE7)--(PTE8) now say exactly that the two forms in
(PTE10) are positive semidefinite.  This proves necessity without
discarding zero eigenvalues.  Conversely, (PTE10) makes \(H\) a direct
sum of a positive semidefinite block and a negative semidefinite block.
Its positive trace is the trace of the first block, which is \(L\) by
(PTE3).  Existence of such metrics is proved constructively below, so
the lower bound is a minimum.

In spectral language, equality in the projection bound means that
\(P_0\) contains every strictly positive eigendirection of \(H\) and
annihilates every strictly negative eigendirection.  It may select a
subspace of \(\ker H\).  The proof above shows that this freedom on the
zero spectral subspace still forces the original positive and negative
primary halves to be metric-orthogonal; it does not force strict
definiteness of the weight on either half.

\subsection{All minimizing metrics as finite Lyapunov expressions}

For either half let \(D\) denote \(B_+\) or \(C_-\), acting on a
\(d\)-dimensional complete space.  Every metric in that half is exactly
\begin{equation}\tag{PTE11}\label{eq:pte11}
 G_D=\int_0^\infty e^{-tD^*}Q_D e^{-tD}\,dt,
 \quad Q_D\succeq0,
 \quad \bigcap_{j=0}^{d-1}\ker(Q_D^{1/2}D^j)=\{0\}.
\end{equation}
To establish this parametrization, write \(D\) in complete primary
coordinates as blocks with positive-real-part eigenvalues.  Each
entry of its exponential is a decaying exponential times a polynomial
of degree at most \(e-1\).  This proves convergence of the integral
entry by entry, including after multiplication by any fixed vector.
Differentiating the integrand and integrating at the two endpoints
gives
\begin{equation}\tag{PTE12}\label{eq:pte12}
 D^*G_D+G_DD=Q_D.
\end{equation}
If a vector \(v\) has zero norm in (PTE11), the nonnegative continuous
integrand \(\|Q_D^{1/2}e^{-tD}v\|^2\) has integral zero and hence
vanishes at every \(t\geq0\).  Its vector-valued analytic square root
has all Taylor coefficients zero, giving
\(Q_D^{1/2}D^jv=0\) for every \(j\geq0\).  Conversely those equations
force the exponential to vanish.  The Cayley--Hamilton identity
expresses every power of degree at least \(d\) through lower powers,
so the first \(d\) equations are equivalent to all of them.  This
proves the positive-definiteness criterion in (PTE11).
For the converse, start with \(G_D>0\) and
\(Q_D=D^*G_D+G_DD\succeq0\).  Integration of
\[
 \frac{d}{dt}(e^{-tD^*}G_De^{-tD})=-e^{-tD^*}Q_De^{-tD}
\]
recovers (PTE11), since the term at infinity is zero by the same
exponential-polynomial decay.  The positive-definiteness criterion
then follows from the argument just given.  This proves both
directions, including singular choices of \(Q_D\).

Use the complete divided-derivative primary frame, ordered with
indices \(0,\ldots,e-1\), and let \(J_e\) have ones on its first
subdiagonal.  It is multiplication by \(S-\lambda\) in this frame.
The two choices of \(D\) are explicitly
\begin{equation}\tag{PTE13}\label{eq:pte13}
 \begin{array}{lll}
 D=B_+:& b_\alpha=\lambda_\alpha-k/2,&N_\alpha=+J_e,\\
 D=C_-:& b_\alpha=k/2-\lambda_\alpha,&N_\alpha=-J_e,
 \end{array}
 \qquad D_\alpha=b_\alpha I+N_\alpha.
\end{equation}
For \(Q_D\) with all original cross-primary blocks retained, the
integral (PTE11) evaluates to
\begin{equation}\tag{PTE14}\label{eq:pte14}
 (G_D)_{\alpha\beta}
 =\sum_{r,s=0}^{e-1}
 \frac{(-1)^{r+s}(r+s)!}
 {r!s!(\overline b_\alpha+b_\beta)^{r+s+1}}
 (N_\alpha^*)^r(Q_D)_{\alpha\beta}N_\beta^s.
\end{equation}
Indeed, expand both finite nilpotent exponentials and integrate
\(t^{r+s}e^{-zt}\).  For \(\operatorname{Re}z>0\), integration by
parts \(r+s\) times gives
\(\int_0^\infty t^{r+s}e^{-zt}dt=(r+s)!/z^{r+s+1}\), with the
boundary terms vanishing.  This proves the formula for complex
denominators as well.  On the negative half, substituting
\(N_\alpha=N_\beta=-J_e\) cancels the factor \((-1)^{r+s}\).
Keeping \(+J_e\) on that half would give incorrect finite entries.
No denominator vanishes, because
\(\operatorname{Re}(\overline b_\alpha+b_\beta)\geq2\delta\).
Conjugation by the original confluent primary-coordinate map returns
these forms to the original polynomial basis.

\subsection{A full-packet minimizer with exactly two nonzero weight eigenvalues}

The preceding parametrization gives a stronger construction than a
strictly definite weight on each primary half.  In each half define
the single linear functional
\begin{equation}\tag{PTE15}\label{eq:pte15}
 \ell_D(v)=\sum_{\alpha}(v_\alpha)_{e-1},
 \qquad Q_D=\ell_D^*\ell_D.
\end{equation}
Every original block appears in this sum.  This \(Q_D\) has rank one
and satisfies the positive-definiteness test in (PTE11).  Here is a
complete verification.  Write \(N_\alpha=\sigma J_e\), with
\(\sigma=+1\) for \(B_+\) and \(\sigma=-1\) for \(C_-\).  If
\(\ell_De^{-tD}v=0\) for every \(t\geq0\), then
\begin{equation}\tag{PTE16}\label{eq:pte16}
 \sum_\alpha e^{-b_\alpha t}p_\alpha(t)=0,
 \qquad
 p_\alpha(t)=\sum_{r=0}^{e-1}
 \frac{(-\sigma t)^r}{r!}(v_\alpha)_{e-1-r}.
\end{equation}
The numbers \(b_\alpha\) are pairwise distinct: equality of real and
imaginary parts in (PTE1), with \(\delta,\gamma>0\), forces the two
indices to agree.  Fix \(\alpha\) and apply
\(\prod_{\beta\ne\alpha}(\partial_t+b_\beta)^e\) to (PTE16).
All terms except \(\alpha\) vanish, leaving
\[
 e^{-b_\alpha t}
 \prod_{\beta\ne\alpha}(\partial_t+b_\beta-b_\alpha)^e
 p_\alpha(t)=0.
\]
On polynomials of degree less than \(e\), each factor
\(\partial_t+c\), \(c\ne0\), is invertible: its inverse is the finite
sum \(c^{-1}\sum_{j=0}^{e-1}(-\partial_t/c)^j\).
Thus \(p_\alpha=0\).  Every coefficient in (PTE16) then gives one
original coordinate of \(v_\alpha\) equal to zero.  Repeating this for
all \(\alpha\) proves \(v=0\), and therefore the finite criterion in
(PTE11).  This argument proves observability of the full generalized
primary chains, not only of their eigenvectors.

In this construction the entries of the positive definite metric
have the following completely explicit form.  For
\(0\leq i,j\leq e-1\), put \(r=e-1-i\) and \(s=e-1-j\).  Then
\begin{equation}\tag{PTE17}\label{eq:pte17}
 (G_D)_{\alpha i,\beta j}
 =\frac{(-\sigma)^{r+s}(r+s)!}
 {r!s!(\overline b_\alpha+b_\beta)^{r+s+1}}.
\end{equation}
Formula (PTE17) follows from (PTE14), since
\((Q_D)_{\alpha\beta}=e_{e-1}e_{e-1}^*\), and the sole contributing
nilpotent powers for this entry are \(r=e-1-i\), \(s=e-1-j\).
In particular it retains every cross term between distinct primary
blocks in the same spectral half.

Take the direct sum of these two positive definite metrics and pull
it back by the original primary-coordinate map.  Call the resulting
original-coordinate metric \(G^{(2)}\).  Its Hermitian weight has
exact spectrum
\begin{equation}\tag{PTE18}\label{eq:pte18}
 \operatorname{spec}H_{G^{(2)}}
 =\{L,-L,\underbrace{0,\ldots,0}_{q-2}\},
 \qquad \operatorname{Tr}(H_{G^{(2)}})_+=L,
 \qquad \|H_{G^{(2)}}\|=L.
\end{equation}
To prove this, the positive-half weight is congruent to
\(Q_{B_+}\), so it is positive semidefinite of rank one.  Its trace
is \(L\) by (PTE3), hence its only nonzero eigenvalue is \(L\).
The negative-half weight is the negative of a positive semidefinite
rank-one matrix, with trace \(-L\).  Its only nonzero eigenvalue is
therefore \(-L\).  This proves (PTE18) and attainment of the minimum.
The zero weight eigenspace has not removed any vector from \(E\):
\(G^{(2)}\) is positive definite on all \(q\) original dimensions,
and \(A\) retains every original Jordan chain.  Any weight attaining
the minimum has at least one positive and at least one negative
eigenvalue, since its positive trace is \(L>0\) and its total trace is
zero.  Thus the rank two in (PTE18) is the smallest possible rank.

More generally every pair of ranks
\begin{equation}\tag{PTE19}\label{eq:pte19}
 1\leq r_+\leq d,\qquad 1\leq r_-\leq d
\end{equation}
occurs among minimizing weights.  Choose a rank-\(r_+\) positive
semidefinite \(Q_{B_+}\) whose range contains \(\ell_{B_+}^*\), and a
rank-\(r_-\) positive semidefinite \(Q_{C_-}\) whose range contains
\(\ell_{C_-}^*\).  For a positive semidefinite matrix the kernel is the
orthogonal complement of its range.  Hence its finite observability
kernel is contained in the one proved zero for \(\ell_D\).  Formula
(PTE11) gives positive definite metrics, and invertible congruence
preserves the ranks of their half-weights.  This proves (PTE19).

For all metrics whose full Hermitian weight has at most one positive
eigenvalue, (PTE8) implies \(\|H_G\|\geq L\).  Construction (PTE18)
attains this bound with one positive and one negative eigenvalue.
This is an exact comparison within the rank-two weight class.  It
does not evaluate the metric selected by the original arithmetic
minimum section; that section has its own fixed Gram matrix.  The
typed section lift between the two metric presentations is the
explicit original-kernel correction treated in the receiving proof.

\subsection{Quantitative excess and the original metric cross term}

The following estimate gives a direct relation between the excess and
the off-diagonal matrix in the exact frame (PTE6):
\begin{equation}\tag{PTE20}\label{eq:pte20}
 \mathcal E(G)\geq\sqrt{L^2+\|C\|_1^2}-L,
 \qquad
 \|C\|_1^2\leq\mathcal E(G)(2L+\mathcal E(G)).
\end{equation}
Here \(\|C\|_1\) is the sum of singular values.  To prove it, write an
ordinary singular value decomposition \(C=U\Sigma V^*\), including
orthonormal bases in the nullspaces, and put \(W=UV^*\).  Then
\[
 T_0=\begin{pmatrix}0&W\\W^*&0\end{pmatrix},
 \qquad T_0^2=J_0^2=I,
 \qquad T_0J_0+J_0T_0=0.
\]
For every real \(\theta\),
\(Z_\theta=\cos\theta J_0+\sin\theta T_0\) is Hermitian with
\(Z_\theta^2=I\).  Diagonalizing \(H\) proves
\(\operatorname{Tr}(Z_\theta H)\leq\operatorname{Tr}|H|\), since
each diagonal entry of a Hermitian contraction lies in \([-1,1]\).
The left side is
\(2L\cos\theta+2\|C\|_1\sin\theta\).  Its maximum is
\(2\sqrt{L^2+\|C\|_1^2}\); the right side is
\(2(L+\mathcal E(G))\).  Squaring this nonnegative inequality proves
both statements in (PTE20).

This inequality is sharp as an inequality for Hermitian block matrices
with these two block traces.  In fact for
\(H=\left(\begin{smallmatrix}aI_d&cW\\cW^*&-aI_d\end{smallmatrix}\right)\),
where \(a>0\), \(c\geq0\), and \(W\) is unitary, one has
\(L=da\), \(\|C\|_1=dc\), and
\(\mathcal E=d(\sqrt{a^2+c^2}-a)\), giving equality.  This last family
establishes sharpness of the block inequality; it does not assert
that all its members have the fixed packet eigenvalues (PTE1).

A second estimate, sometimes stronger, is
\begin{equation}\tag{PTE21}\label{eq:pte21}
 \|C\|_{\mathrm F}^2
 \leq\bigl(\|P_0H_+P_0\|+\|Q_0H_-Q_0\|\bigr)\mathcal E(G)
 \leq\bigl(\|H_+\|+\|H_-\|\bigr)\mathcal E(G).
\end{equation}
For a positive semidefinite matrix \(T\),
\[
 \|P_0TQ_0\|_{\mathrm F}^2
 =\operatorname{Tr}\bigl[(T^{1/2}P_0T^{1/2})
                         (T^{1/2}Q_0T^{1/2})\bigr]
 \leq\|P_0TP_0\|\operatorname{Tr}(Q_0T).
\]
The norm equality used here follows because \(R^*R\) and \(RR^*\)
have the same largest eigenvalue, with \(R=P_0T^{1/2}\).
Apply this inequality to \(T=H_+\), and apply its exchanged-projection
version to \(T=H_-\).  Since
\(C=P_0H_+Q_0-P_0H_-Q_0\), the triangle inequality followed by the
two-term scalar Cauchy inequality bounds its Frobenius norm by
\[
 \sqrt{\|P_0H_+P_0\|\operatorname{Tr}(Q_0H_+)}
 +\sqrt{\|Q_0H_-Q_0\|\operatorname{Tr}(P_0H_-)}.
\]
Squaring and using (PTE8) gives (PTE21).

The excess also has the exact decomposition
\begin{equation}\tag{PTE22}\label{eq:pte22}
 \mathcal E(G)=\operatorname{Tr}(D_+)_-
 +\operatorname{Tr}(D_-)_+
 +\frac12\bigl(\|H\|_1-\|D_+\|_1-\|D_-\|_1\bigr),
 \qquad \|H\|_1-\|D_+\|_1-\|D_-\|_1\geq0.
\end{equation}
The identity follows by substituting the traces in (PTE7).
For the inequality, choose the block-diagonal Hermitian contraction
\(\operatorname{diag}(\operatorname{sgn}D_+,
\operatorname{sgn}D_-)\); its trace pairing with \(H\) is
\(\|D_+\|_1+\|D_-\|_1\), bounded above by \(\|H\|_1\) by the
same diagonalization argument as above.  Equation (PTE22) therefore
retains both within-half sign defects as well as the cross-block
contribution.

There is also an explicit inverse from the measured cross block
\(C\) back to the original metric cross term \(X\).  Write the
blocks of \(B_+\) as \(b_\alpha I+J_e\), and those of \(B_-\) as
\(\beta_\eta I+J_e\), so
\(\operatorname{Re}b_\alpha\geq\delta\) and
\(\operatorname{Re}\beta_\eta\leq-\delta\).  Define
\(R=P^{-1/2}CS^{1/2}=KB_--B_+K\).  Then
\begin{equation}\tag{PTE23}\label{eq:pte23}
 K_{\alpha\eta}
 =-\sum_{r,s=0}^{e-1}
 \frac{(-1)^r(r+s)!}
 {r!s!(b_\alpha-\beta_\eta)^{r+s+1}}
 J_e^rR_{\alpha\eta}J_e^s,
 \qquad X=PK.
\end{equation}
Indeed the convergent integral
\[
 K=-\int_0^\infty e^{-tB_+}R e^{tB_-}\,dt
\]
solves \(B_+K-KB_-=-R\), by differentiation of the integrand and
evaluation of its endpoints.  In complete primary coordinates its
expansion is (PTE23).  For uniqueness, a solution of the homogeneous
equation satisfies
\(e^{-tB_+}Ke^{tB_-}=K\) for all \(t\); the left side tends to zero
because both exponentials decay with finite polynomial factors.
Thus \(K=0\).  This also supplies the exact invertibility invoked in
the equality proof.

For an explicit norm version put
\begin{equation}\tag{PTE24}\label{eq:pte24}
 \kappa=\max_{\alpha,\eta}
 \sum_{r,s=0}^{e-1}
 \frac{(r+s)!}{r!s!\,|b_\alpha-\beta_\eta|^{r+s+1}}.
\end{equation}
Every nilpotent power of the shift has operator norm at most one.
Apply the Frobenius triangle inequality to each block of (PTE23),
square and sum over blocks, and use
\(\|R\|_{\mathrm F}\leq\|P^{-1/2}\|\|S^{1/2}\|\|C\|_{\mathrm F}\).
Together with (PTE20), this proves
\begin{equation}\tag{PTE25}\label{eq:pte25}
 \|P^{-1}X\|_{\mathrm F}^2
 \leq\kappa^2\|P^{-1/2}\|^2\|S^{1/2}\|^2
 \mathcal E(G)(2L+\mathcal E(G)).
\end{equation}
The original dependence on the metric restrictions and its Schur
complement is displayed explicitly.  No uniform estimate in those
metric factors has been assumed.

\subsection{Exact primary-coordinate polynomial matrix identities}

The sign and rank assertions can be tested by a finite calculation
without substituting approximate roots.  For each spectral half take
distinct positive-real-part complex numbers \(b_\alpha\), retain
\(N_\alpha=\sigma J_e\), and form the matrix in (PTE17).  Equations
(PTE12), (PTE15), and (PTE17) imply entry by entry
\begin{equation}\tag{PTE26}\label{eq:pte26}
 (\overline b_\alpha+b_\beta)G_{\alpha i,\beta j}
 +\sigma\,\mathbf 1_{i<e-1}G_{\alpha,i+1;\beta j}
 +\sigma\,\mathbf 1_{j<e-1}G_{\alpha i;\beta,j+1}
 =\mathbf 1_{i=e-1}\mathbf 1_{j=e-1}.
\end{equation}
For \(r+s>0\), substitute (PTE17).  The common nonzero denominator
is \((\overline b_\alpha+b_\beta)^{r+s}\), and the three numerators
reduce to
\[
 (-\sigma)^{r+s}
 \left[\frac{(r+s)!}{r!s!}
 -\mathbf1_{r>0}\frac{(r+s-1)!}{(r-1)!s!}
 -\mathbf1_{s>0}\frac{(r+s-1)!}{r!(s-1)!}\right]=0.
\]
For \(r=s=0\) only the first term remains and equals one.
This proves (PTE26) directly, including the boundary rows and columns,
and verifies the negative-half signs independently of integration.

\paragraph{Source and precise scope.}
The received source is \texttt{04\_FULL\_METRIC\_POSITIVE\_TRACE\_MINIMUM.md},
19 September 2026, formulas (M1)--(M7).  Its minimum and equality
classification are confirmed by the complete proof above.  The explicit
negative-half sign in (PTE13)--(PTE14), rank-two construction
(PTE15)--(PTE19), excess identities (PTE8)--(PTE9), and original
cross-metric estimates (PTE20)--(PTE25) are the independent continuation.
All algebra used in this fragment is proved here.  The results optimize
the full-packet Hermitian positive trace and characterize its attaining
metrics.  The original arithmetic attained quotient, its seven absolute
allocations, and its signed projected currents still require their
own evaluated Gram and projection data; no value for those quantities
is assigned by this finite-dimensional optimization.
